\documentclass[11pt,a4paper]{article}
\usepackage[margin=27mm]{geometry}
\usepackage{amsmath,amssymb,amsthm}
\usepackage[T1]{fontenc}
\usepackage{lmodern}
\usepackage{microtype}
\usepackage[hidelinks]{hyperref}
\newcommand{\R}{\mathbb R}
\newcommand{\dd}{\,d}
\DeclareMathOperator{\Tr}{Tr}
\DeclareMathOperator{\Per}{Per}
\newtheorem{theorem}{Theorem}[section]
\newtheorem{lemma}[theorem]{Lemma}
\newtheorem{proposition}[theorem]{Proposition}
\theoremstyle{remark}
\newtheorem{remark}[theorem]{Remark}
\numberwithin{equation}{section}
\title{Concentric balls for thermal insulation\\with a radiative boundary law}
\author{Kaiwen Jin\\[2pt]\small\href{mailto:yui.ui.siki@gmail.com}{yui.ui.siki@gmail.com}}
\date{15 October 2026}
\hypersetup{pdftitle={Concentric balls for thermal insulation with a radiative boundary law},pdfauthor={Kaiwen Jin}}
\begin{document}
\maketitle
\begin{abstract}
We consider a thermal insulation energy with a quadratic Dirichlet term
and the nonlinear boundary potential associated with radiative heat
transfer. Both the heated set and the surrounding domain may vary,
with heated volume normalized to that of the unit ball and an upper
bound on the total volume. We prove that the minimum is achieved within the family of
concentric balls, allowing a smaller outer radius and the bare heated
ball. The comparison holds in every dimension \(n\ge2\), for every
positive boundary coefficient and every nonnegative ambient parameter.
The proof rearranges the boundary potential and retains its singular
variation. A concavity property of the radiative law then removes
the resulting interior jumps without increasing the energy.
\end{abstract}
\section{Introduction}\label{sec:introduction}

In a shape optimization problem for thermal insulation, one may choose
both the heated body and the region surrounding it. The question is
whether a spherical body with a concentric insulating layer uses a
given amount of material most efficiently. For a stationary model
with quadratic bulk energy, the boundary loss is described by a
nonnegative potential of the temperature. The choice of that
potential distinguishes convective and radiative boundary laws
\cite{BNNT}.

Bucur, Nahon, Nitsch and Trombetti \cite{BNNT} proved concentric-ball
optimality for the quadratic boundary potential. They also obtained
general results for a volume penalty and certain constrained cases
with inactive volume constraint. Their question after Remark~4 asks
whether the concentric comparison holds for more general boundary
potentials, including the radiative law in their equation~(3).
The penalized comparison does not by itself give the constrained
comparison at every prescribed volume. We prove that comparison
for the full radiative potential, with all its parameters retained.

Here is the precise problem. Fix an integer \(n\ge2\), numbers
\(\beta>0\), \(\gamma\ge0\), and \(R>1\), and set
\begin{equation}\label{eq:theta}
 \Theta(s)=\beta\left(\frac{s^5}{5}+\gamma s^4+
                                  2\gamma^2s^3+2\gamma^3s^2\right).
\end{equation}
For a bounded Lipschitz open set \(\Omega\subset\R^n\) and a measurable
set \(K\subset\Omega\), let
\begin{equation}\label{eq:energy}
 E_\Theta(K,\Omega)=
 \min_{\substack{v\in H^1(\Omega),\ 0\le v\le1\\
                         v=1\ {\rm a.e.\ on}\ K}}
 \left\{\int_\Omega|\nabla v|^2\dd x+
          \int_{\partial\Omega}\Theta(\Tr v)\dd\mathcal H^{n-1}\right\}.
\end{equation}
The interval bounds are understood almost everywhere, and \(\Tr v\)
is the interior trace. Write the expression in braces as
\(\mathcal E_\Omega(v)\). We use the unit ball \(B_1\) to normalize
the heated volume and write \(\omega_n=|B_1|\).

\begin{theorem}\label{thm:main}
For every choice of \(n\ge2,\beta>0,\gamma\ge0,R>1\) above, every bounded
Lipschitz open \(\Omega\subset\R^n\), and every measurable \(K\subset\Omega\)
satisfying
\[
 |K|=\omega_n,\qquad |\Omega|\le\omega_nR^n,
\]
one has
\begin{equation}\label{eq:main}
 E_\Theta(K,\Omega)\ge\min_{1\le r\le R}E_\Theta(B_1,B_r).
\end{equation}
At \(r=1\), the pair is \(K=\Omega=B_1\).
\end{theorem}

Here \(B_r\) denotes the ball of radius \(r\) with the same center as
\(B_1\), and the energy on the right is the same trace energy
\eqref{eq:energy}. No connectedness of \(\Omega\) or regularity of
\(K\) is required. The minimum includes the bare ball and can use
less than the full available outer volume. This is the radiative
case of the constrained concentric-ball question, with the
upper-volume convention used in \cite{BNNT}.
We make no assertion here for an arbitrary increasing boundary potential.

The proof uses the boundary potential itself as the rearranged
variable. Extending that potential by zero records the outer
boundary loss as variation. Isoperimetry then gives a radial lower
cost, but the radial profiles may have jumps and need not represent
admissible \(H^1\) temperatures. We approximate their cost by profiles
with finitely many jumps. For the radiative polynomial, multiplying
the potential on the entire tail after a jump has a concave total
cost. One endpoint of the allowed multiplication interval removes
the jump, or removes the whole tail, without raising the energy.
Repeating this operation leaves a genuine concentric-ball competitor.
The possibility of shortening the tail explains why smaller outer
radii, including the bare ball, belong in \eqref{eq:main}.

Section~\ref{sec:radial} gives the radial energy formula and the scalar
concavity property. Section~\ref{sec:rearrangement} proves the radial
lower comparison with all singular variation retained.
Section~\ref{sec:jumps} constructs the finite-jump approximation,
removes the jumps, and proves the theorem.
All functions are real, all volume integrals use Lebesgue measure,
and Hausdorff measure is normalized by
\(\mathcal H^{n-1}(\partial B_r)=n\omega_n r^{n-1}\).

\section{Radial energies and the radiation potential}\label{sec:radial}

Throughout the proof the parameters \(n,\beta,\gamma,R\) are fixed as in
Theorem~\ref{thm:main}. We write
\[
 S=n\omega_n,\qquad A(r)=r^{n-1},\qquad Y=\Theta(1).
\]
Thus \(S A(r)\) is the area of the sphere of radius \(r\).

The interior trace on a bounded Lipschitz open set is a bounded linear
map from \(BV\) into boundary \(L^1\). With outward normal \(\nu\),
Gauss--Green reads
\[
 \int_\Omega f\,\operatorname{div}X+\int_\Omega X\cdot\dd Df
 =\int_{\partial\Omega}\Tr f\,X\cdot\nu\,\dd\mathcal H^{n-1}.
\]
This is the inner-normal formula of \cite[Theorem~2.9]{LV} with its sign
converted. We use the absolute-mean characterization
\[
 \lim_{\rho\downarrow0}\frac1{|B_\rho(x)\cap\Omega|}
 \int_{B_\rho(x)\cap\Omega}|f-\Tr f(x)|\dd y=0
 \quad\text{for }\mathcal H^{n-1}\text{-a.e. }x\in\partial\Omega
\]
from \cite[Corollary~1.29]{Gorny}.\footnote{The proof of this
corollary is omitted in the lecture notes, whose publication date is
not specified. The point in \cite[(2.5)]{LV} is printed as \(x\in\Omega\),
whereas \cite[(2.6)]{LV} is written on \(\partial\Omega\). The
absolute-mean statement used here is the boundary statement in
\cite[Corollary~1.29]{Gorny}.}
If a connected domain is required in that characterization, apply it
in one-sided Lipschitz chart cylinders. On smaller boundary patches the
small-ball intersections agree with those in \(\Omega\), and uniqueness
of the local trace gives the stated formula without connectedness of
\(\Omega\).


We first justify the minimum in the definition of the energy. An admissible
minimizing sequence is bounded in \(H^1(\Omega)\), since its values lie in
\([0,1]\). Pass to a weakly convergent subsequence. Tests by indicators show
that the limit still lies in \([0,1]\) and equals 1 almost everywhere on \(K\).
The bounded trace map gives
\[
 \|\Tr f\|_{L^1(\partial\Omega)}
 \le C_\Omega\bigl(\|f\|_{L^1(\Omega)}+\|\nabla f\|_{L^1(\Omega)}\bigr).
\]
Consequently every pairing of the trace with a bounded boundary function
is a continuous linear functional on \(H^1(\Omega)\). Traces of admissible
temperatures also lie in \([0,1]\). Since
\(\Theta''(s)=4\beta(s+\gamma)^3\ge0\), the supporting-line inequality at
the limiting trace proves weak lower semicontinuity of the boundary
integral. The gradient integral has the same property. The constant
temperature 1 is an admissible finite-energy test, so the minimum is attained.

\begin{lemma}\label{lem:radial}
For \(r>1\), let \(I_n(r)=\int_1^r \rho^{1-n}\dd\rho\). Then
\begin{equation}\label{eq:radial}
 \frac{E_\Theta(B_1,B_r)}{S}
 =\min_{0\le t\le1}\left\{\frac{(1-t)^2}{I_n(r)}+
                                      r^{n-1}\Theta(t)\right\}.
\end{equation}
The energy is continuous on \([1,R]\) when its value at 1 is defined
by the bare ball, \(E_\Theta(B_1,B_1)=S Y\). In particular its minimum
on this interval is attained.
\end{lemma}

\begin{proof}
For an admissible \(v\) on \(B_r\), let
\[
 m(\rho)=\frac1S\int_{\partial B_1}v(\rho\theta)
                                      \dd\mathcal H^{n-1}(\theta).
\]
Radial tests in the weak-derivative identity give
\(m'(\rho)=S^{-1}\int_{\partial B_1}\nabla v(\rho\theta)\cdot\theta
\dd\mathcal H^{n-1}(\theta)\). Cauchy--Schwarz therefore yields
\[
 S\int_1^r A(\rho)|m'(\rho)|^2\dd\rho
 \le\int_{B_r\setminus B_1}|\nabla v|^2\dd x.
\]
In particular \(m\) has an absolutely continuous representative on
\([1,r]\). Applying the same argument across radius 1 shows \(m(1)=1\).
Its value \(m(r)\) is the mean of the actual boundary trace. To see this
without selecting a boundary slice, use Gauss--Green with the radial
field \(f(\rho)\theta\), zero near the origin. Polar integration turns
its left side into
\[
 S\int_0^r\bigl[m(Af)'+m'Af\bigr]\dd\rho
       =S A(r)f(r)m(r),
\]
whereas its boundary side is
\(f(r)\int_{\partial B_r}\Tr v\,\dd\mathcal H^{n-1}\).
Here \(m'=0\) in the core.

Put \(t=m(r)\in[0,1]\). Weighted Cauchy--Schwarz gives
\((1-t)^2\le I_n(r)\int_1^r A|m'|^2\).
The supporting line for \(\Theta\) at the boundary mean gives
\(\int_{\partial B_r}\Theta(\Tr v)\ge S A(r)\Theta(t)\).
Conversely the radial temperature
\[
 u(\rho)=1-(1-t)\frac{I_n(\rho)}{I_n(r)},\qquad 1\le\rho\le r,
\]
extended by 1 in the core, has exactly the cost on the right of
\eqref{eq:radial}. This proves the formula.

Compact minimization in \(t\) gives continuity for \(r>1\).
For a minimizing \(t_r\), the test \(t=1\) implies
\[
 (1-t_r)^2\le I_n(r)r^{n-1}Y\longrightarrow0\qquad(r\downarrow1).
\]
Thus its boundary term tends to \(Y\), while the entire expression
is at most \(r^{n-1}Y\to Y\). This proves continuity at 1 and the
assertion about the minimum. The integral definition includes
\(I_2(r)=\log r\), so no separate singular formula occurs in dimension two.
\end{proof}

We next identify the scalar property used to remove jumps. The potential
\(\Theta\) is strictly increasing on \([0,1]\), because
\(\Theta'(s)>0\) for \(s>0\). Let \(Q:[0,Y]\to[0,1]\) be its continuous
inverse, and define
\begin{equation}\label{eq:GH}
 G(y)=\frac{\Theta'(Q(y))^2}{4},\qquad
 H(y)=\frac{y^2}{G(y)}\quad(y>0),\qquad G(0)=H(0)=0.
\end{equation}
In potential variables \(y=\Theta(u)\), the Dirichlet density is
\(|y'|^2/(4G(y))\). The function \(H\) appears when this potential
is multiplied by a constant.

\begin{lemma}\label{lem:concavity}
The function \(G\) is Lipschitz on \([0,Y]\), positive on \((0,Y]\),
and \(H\) is concave on \([0,Y]\).
\end{lemma}

\begin{proof}
For \(y>0\), inverse differentiation gives
\[
 G'(y)=\frac{\Theta''(Q(y))}{2}
       =2\beta(Q(y)+\gamma)^3.
\]
This is bounded by \(2\beta(1+\gamma)^3\), and \(G\) is continuous at
zero. The Lipschitz assertion follows.
If \(\gamma=0\), then
\[
 H(y)=\frac4{25}\left(\frac{5y}{\beta}\right)^{2/5},
\]
which has the required concavity. Suppose \(\gamma>0\), put \(z=s/\gamma\),
and set
\[
 P(z)=z^5/5+z^4+2z^3+2z^2.
\]
Then \(\Theta(s)=\beta\gamma^5P(z)\) and
\(H(\Theta(s))=4\gamma^2P(z)^2/P'(z)^2\).
For \(h=4P^2/P'^2\), differentiation gives
\[
 h'=8P/P'-8P^2P''/P'^3,\qquad
 h''=8-24PP''/P'^2-8P^2P'''/P'^3+24P^2P''^2/P'^4,
\]
and hence
\begin{equation}\label{eq:concavitynumerator}
 h''P'-h'P''=-\frac{8}{P'^3}
          \bigl[P^2P'P'''-(P'^2-2PP'')^2\bigr].
\end{equation}
The second derivative of \(h\) with respect to \(P\) is the left
side of \eqref{eq:concavitynumerator} divided by \(P'^3\).
For \(q=5P=z^5+5z^4+10z^3+10z^2\), direct multiplication gives
\begin{align}
 q^2q'q'''-(q'^2-2qq'')^2
 =25z^5\bigl(&3z^{11}+48z^{10}+360z^9+1680z^8+5508z^7 \notag\\
 &+13584z^6+26136z^5+39600z^4+46080z^3 \notag\\
 &+38400z^2+20000z+4800\bigr). \label{eq:positivepoly}
\end{align}
All coefficients are positive. The expression in brackets in
\eqref{eq:concavitynumerator} is \(1/625\) times the left side of
\eqref{eq:positivepoly}. Since \(P'>0\), this proves strict concavity
at positive potential values; all changes of scale have positive factors.
Finally \(P(z)\sim2z^2\), \(P'(z)\sim4z\), so
\(H(\Theta(s))\sim s^2\to0\) as \(s\downarrow0\).
The continuous zero extension preserves concavity.
\end{proof}

\section{Rearranging the boundary potential}\label{sec:rearrangement}

The boundary term becomes linear in the variable \(U=\Theta(v)\).
After zero extension this potential is a \(BV\) function, and may have
a boundary jump. Rearrangement must therefore retain singular variation.
We recall the formulas used below, with
\(\Per(E)=|D\mathbf1_E|(\R^n)\).
For nonnegative \(V\in BV(\R^n)\) of bounded support, coarea gives
\[
 |DV|(\R^n)=\int_0^\infty\Per(\{V>t\})\dd t;
\]
the converse holds when the right side is finite and \(V\in L^1\).
We use its Borel-set version as well
\cite[Theorem~1.19]{Gorny}; see also
\cite[Theorem~2.6]{LV}.
Euclidean isoperimetry in this normalization is
\[
 \Per(E)\ge n\omega_n^{1/n}|E|^{(n-1)/n}
\]
for bounded measurable sets, with equality for balls
\cite[Theorem~2.1]{AGNT}.

For \(f\in W^{1,1}(\Omega)\), Gauss--Green gives the zero-extension
identity
\begin{equation}\label{eq:zeroextension}
 D(f\mathbf1_\Omega)=\nabla f\,\dd x|_\Omega-
                     \Tr f\,\nu\,\mathcal H^{n-1}|_{\partial\Omega},
 \qquad
 |D(f\mathbf1_\Omega)|(\R^n)=\int_\Omega|\nabla f|+
                                          \int_{\partial\Omega}|\Tr f|.
\end{equation}
The volume and boundary measures are mutually singular.
The absolute-mean characterization also shows that the trace commutes
with a Lipschitz scalar composition.

For a Borel \(p:[0,Y]\to[0,1]\), let \(B(t)=\int_0^t p(a)\dd a\).
This is a nondecreasing 1-Lipschitz function. The Sobolev chain rule gives
\(\nabla B(U)=p(U)\nabla U\) inside \(\Omega\); the \(C^1\) version is
\cite[Proposition~3.21(2)]{Hunter}, and we spell out the representative
issue. Sobolev functions have absolutely continuous representatives on
almost every coordinate line: product tests and Fubini give the
one-dimensional weak derivative, whose integral differs from the
function by a constant \cite[Proposition~3.6]{Hunter}.
For an absolutely continuous real \(f\) and a real null Borel set \(N\),
\(f'=0\) a.e. on \(f^{-1}(N)\). Indeed, points with nonzero derivative
can be grouped into countably many sets on which the inverse
restriction of \(f\) is Lipschitz: group by integers \(k,l\) for which
\(|f(y)-f(x)|\ge |y-x|/k\) when \(0<|y-x|<1/l\), and then by intervals
of length less than \(1/l\). Their points over \(N\) are null, since
a Lipschitz map sends null sets to null sets. Apply this to the
exceptional levels of \(B'\). The ordinary chain rule holds elsewhere,
and the Lipschitz bound gives derivative zero where \(f'=0\).
Fubini proves the asserted weak chain rule with the chosen Borel \(p\).

We will also use the weighted coarea identity
\begin{equation}\label{eq:weightedcoarea}
 |D(B(V))|(\R^n)=\int_0^Y p(t)\Per(\{V>t\})\dd t
\end{equation}
for \(0\le V\le Y\) in \(BV\) of bounded support. To include flat pieces
of \(B\), put
\(\tau(q)=\sup\{t\in[0,Y]:B(t)\le q\}\) for \(0<q<B(Y)\).
Then \(\{B(V)>q\}=\{V>\tau(q)\}\). The one-dimensional area formula
\cite[Theorem~2.8]{LV}, first on indicator functions and then on
nonnegative measurable weights, changes \(dq\) to \(p(t)\dd t\).
Non-singleton fibers of \(B\) have only countably many distinct values.
This proves \eqref{eq:weightedcoarea}; its finite right side first
justifies \(B(V)\in BV\) by the coarea converse. If \(B(Y)=0\), both
sides are zero.

The scalar cost of an absolutely continuous potential drop will be
\begin{equation}\label{eq:psi}
 \psi(y,z)=\sup_{0\le p\le1}\{pz-G(y)p^2\},\qquad0\le y\le Y,\quad z\ge0.
\end{equation}
For \(G(y)>0\) this is
\[
 \psi(y,z)=
 \begin{cases}
 z^2/(4G(y)),&z\le2G(y),\\
 z-G(y),&z\ge2G(y).
 \end{cases}
\]
At \(y=0\) it is \(z\). The supremum formula shows that it is convex,
nondecreasing and 1-Lipschitz in \(z\), that \(0\le\psi(y,z)\le z\),
and that its recession slope is 1. Writing
\(\psi_g(z)=\sup_{0\le p\le1}(pz-gp^2)\), we have
\begin{equation}\label{eq:psilip}
 |\psi_{g_1}(z)-\psi_{g_2}(z)|\le |g_1-g_2| .
\end{equation}

Let \(w\) be a nonincreasing profile on \((1,R)\), taking values in
\([0,Y]\). Augment it by \(Y\) to the left of 1 and by 0 to the right
of \(R\). Its positive drop measure \(\mu=-Dw\) on \([1,R]\) has total
mass \(Y\), including the atoms \(Y-w(1+)\) at 1 and \(w(R-)\) at \(R\).
Decompose it as
\(\mu=z(r)\dd r+\mu_s\), with \(\mu_s\perp\dd r\), and define
\begin{equation}\label{eq:F}
 F(w)=\int_1^R A(r)\psi(w(r),z(r))\dd r+
                              \int_{[1,R]}A(r)\dd\mu_s(r).
\end{equation}
The singular coefficient is the recession slope of \(\psi\).
It includes jump and Cantor variation.

\begin{proposition}\label{prop:lower}
For every original admissible temperature \(v\), there is such a
profile \(w\) for which
\begin{equation}\label{eq:lower}
 \mathcal E_\Omega(v)\ge S F(w).
\end{equation}
\end{proposition}

\begin{proof}
Extend \(U=\Theta(v)\) by zero outside \(\Omega\). The \(C^1\) chain
rule and \eqref{eq:zeroextension} put \(U\) in \(BV(\R^n)\).
Completing a square gives
\[
 |\nabla v|^2\ge p(U)|\nabla U|-G(U)p(U)^2 .
\]
Since \(B(\Theta(\Tr v))\le\Theta(\Tr v)\), we obtain
\begin{equation}\label{eq:dualoriginal}
 \mathcal E_\Omega(v)\ge |D(B(U))|(\R^n)-
                                       \int_{\R^n}G(U)p(U)^2\dd x.
\end{equation}
Let \(U^*(x)=w(|x|)\) be the decreasing spherical rearrangement.
Explicitly, for \(m(t)=|\{U>t\}|\), take
\[
 w(r)=\sup\{t\in[0,Y]:m(t)>\omega_n r^n\},
\]
with empty supremum 0. Right continuity of \(m\) shows that
\(\{U^*>t\}\) is the ball of volume \(m(t)\).
Isoperimetry and the coarea converse give \(U^*\in BV\).
By \eqref{eq:weightedcoarea}, the variation of \(B(U^*)\) does not
exceed that of \(B(U)\); equimeasurability preserves the bulk term
in \eqref{eq:dualoriginal}, since its integrand vanishes at zero.
Moreover \(U=Y\) on \(K\) and its support has volume at most
\(\omega_nR^n\). Thus \(w=Y\) on \(r<1\) and \(w=0\) on \(r>R\), and
\begin{equation}\label{eq:radialdual}
 \mathcal E_\Omega(v)\ge
 \sup_p\left\{|D(B(U^*))|-\int_{\R^n}G(U^*)p(U^*)^2\dd x\right\}.
\end{equation}

We compute this supremum without assuming an absolutely continuous
inverse distribution. Complete the monotone graph of \(w\) by filling
all vertical jumps, including the endpoint jumps. From \((1,Y)\)
to \((R,0)\), use the parameter
\[
 s=r+Y-t,\qquad1\le s\le R+Y.
\]
An increase in \(s\) is exactly the radial increase plus the potential
drop. Hence its coordinate maps \(r(s),t(s)\) are 1-Lipschitz,
and, almost everywhere,
\[
 a=r'(s)\ge0,\qquad b=-t'(s)\ge0,\qquad a+b=1.
\]
There are no missing parameters: the jumps of the strictly increasing
map \(r\mapsto r+Y-w(r)\) are precisely the inserted vertical segments,
and its generalized inverse supplies the rest of the graph.
The one-dimensional area formula gives
\begin{equation}\label{eq:pushforwards}
 r_\#(a\,\dd s)=\dd r,\qquad
 t_\#(b\,\dd s)=\dd t|_{(0,Y)},\qquad r_\#(b\,\dd s)=\mu.
\end{equation}
For the last equality, a radial interval traverses exactly the drop
of \(w\) on it, with its endpoint vertical segments. The area formula
for \(t\) measures that drop. Equality on intervals determines the
finite measures, and nonnegative simple functions give all weights.

On \(a>0\), the radial inverse is unique except at countably many jump
radii; the drop density is \(b/a\). The image of \(a=0\) under \(r\)
is Lebesgue-null. Uniqueness of the Lebesgue decomposition therefore
identifies these as the absolutely continuous and singular parts of
\(\mu\), respectively. At a nonjump radius \(t(s)=w(r(s))\); a vertical
fiber has \(a=0\). Consequently
\begin{equation}\label{eq:graphcost}
 F(w)=\int_1^{R+Y}A(r(s))
       \sup_{0\le p\le1}\{p b(s)-G(t(s))p^2a(s)\}\dd s.
\end{equation}
For \(a>0\) the integrand is \(a\psi(t,b/a)\), and for \(a=0\) it
is \(b\), retaining the whole singular cost.

The radial dual in \eqref{eq:radialdual}, divided by \(S\), equals
\begin{equation}\label{eq:graphdual}
 \int_1^{R+Y}A(r(s))
       [p(t(s))b(s)-G(t(s))p(t(s))^2a(s)]\dd s
                                  -\frac{G(Y)p(Y)^2}{n}.
\end{equation}
The last term is the contribution of the core \(B_1\).
Indeed \eqref{eq:weightedcoarea} integrates the spherical crossing
radius against \(p(t)\dd t\), while polar volume is
\(S A(r)\dd r\); \eqref{eq:pushforwards} gives the two graph integrals.

It remains to choose a single Borel function of the level which
maximizes them. Choose Borel representatives of \(a,b\). A null
Borel set \(N\subset[0,Y]\) can contain \(0,Y\), the image under \(t\)
of \(b=0\), and the image of the parameter-null set where derivatives
are undefined. The area formula and the Lipschitz bound make these
images null. Outside \(N\), \(t\) has a unique inverse parameter
\(s(t)\); otherwise its fiber would contain a horizontal interval.
This generalized inverse is Borel. Set \(p=0\) on \(N\) and elsewhere
choose
\[
 p(t)=\min\{b(s(t))/(2G(t)a(s(t))),1\}\quad\text{if }a(s(t))>0,
 \qquad p(t)=1\quad\text{if }a(s(t))=0.
\]
Here \(G(t)>0\). On \(b=0\), zero is the maximizer. The part of
\(b>0\) with \(t(s)\in N\) has zero parameter measure, since its
strictly positive \(b\)-weighted integral is zero by
\eqref{eq:pushforwards}. Everywhere else the choice is a pointwise
maximizer. It also has \(p(Y)=0\), eliminating the core term.
Conversely every \(p\) gives at most the integral in
\eqref{eq:graphcost}, with a nonpositive core term.
Thus the supremum is exactly \(S F(w)\), proving \eqref{eq:lower}.
\end{proof}

\section{Approximation and removal of jumps}\label{sec:jumps}

The lower bound in Proposition~\ref{prop:lower} involves a monotone
\(BV\) potential. We first approximate its cost by temperatures with
finitely many jumps. Concavity will then remove those jumps.
For a nonincreasing piecewise \(H^1\) temperature \(u\), equal to 1
inside \(B_1\) and zero after radius \(R\), set
\begin{equation}\label{eq:J}
 J(u)=\int_1^R A(r)|u'(r)|^2\dd r+
       \sum_s A(s)\bigl[\Theta(u(s-))-\Theta(u(s+))\bigr].
\end{equation}
The sum runs over its finitely many downward jumps, including a possible
jump at 1 and the terminal jump to zero. This is an auxiliary cost:
interior jumps are not admissible in the original \(H^1\) problem.

\begin{proposition}\label{prop:approximation}
For every augmented monotone profile \(w\) in \eqref{eq:F} and every
\(\epsilon>0\), there is a temperature of the above form satisfying
\[
 J(u)\le F(w)+\epsilon .
\]
\end{proposition}

\begin{proof}
All approximation parameters are chosen with \(n,\beta,\gamma,R\)
fixed. Write \(A_*=R^{n-1}\), \(L_A=(n-1)R^{n-2}\), and
\(L_G=\operatorname{Lip}(G)\).

\emph{Polygonal potentials.}
Choose a partition of \([1,R]\), of mesh \(h\), with all interior
vertices at continuity points of \(w\). Join the endpoint values
linearly on each cell, using \(w(1+)\) and \(w(R-)\) at the ends,
and keep the two endpoint atoms. Call the result \(w_h\).
For a cell \(C_i=(x_i,x_{i+1})\), let \(d_i\) be its length,
\(y_i\) its left value, and \(m_i\ge0\) its total drop. Thus
\(m_i=\mu(C_i)\), with the retained end atoms excluded, and
\(\sum_i m_i\le Y\).

Freeze \(A_i=A(x_i)\) and \(g_i=G(y_i)\).
For every constant \(p\in[0,1]\),
\[
 pm_i-d_ig_ip^2
   =\int_{C_i}(pz(r)-g_ip^2)\dd r+p\mu_s(C_i)
   \le\int_{C_i}\psi_{g_i}(z(r))\dd r+\mu_s(C_i).
\]
Taking the supremum proves
\[
 d_i\psi_{g_i}(m_i/d_i)
       \le\int_{C_i}\psi_{g_i}(z(r))\dd r+\mu_s(C_i).
\]
For both profiles, the potential stays between \(y_i-m_i\) and \(y_i\).
Unfreezing the weight costs at most \(L_Ahm_i\), since
\(\psi\le z\) and the singular coefficient is 1.
By \eqref{eq:psilip}, unfreezing \(G\) costs at most
\(A_*L_Gd_im_i\). Apply these estimates on both sides and sum:
\begin{equation}\label{eq:polygonalerror}
 F(w_h)\le F(w)+2h(L_A+A_*L_G)Y .
\end{equation}
The endpoint atom costs are identical. This estimate does not require
bounded slopes of \(w\), and applies also to Cantor variation.

\emph{Cutoff near zero.}
Fix \(0<\delta<Y\). Follow the polygonal profile until its potential
first falls to or below \(\delta\), and replace the remainder by zero.
At a continuous crossing this adds a terminal drop \(\delta\).
If an existing jump skips the level, its right value is
\(c\le\delta\); replacing that jump by a drop to zero adds only \(c\).
The same reasoning applies at 1 and \(R\). Removed energy is nonnegative.
For the resulting \(q\),
\begin{equation}\label{eq:cutofferror}
 F(q)\le F(w_h)+A_*\delta.
\end{equation}
If the cutoff is immediate at 1, its cost is the bare-ball cost \(Y\).
Otherwise all positive one-sided values of the retained finite polygon
lie in \([\delta,Y]\), up to its terminal jump at some
\(\ell\in(1,R]\).

\emph{Finite jumps and affine segments.}
Keep the existing jumps of \(q\). Let \(D<\infty\) bound its finitely
many affine slopes, and put
\(g_*=\min_{[\delta,Y]}G>0\), \(g^*=\max_{[\delta,Y]}G\).
Subdivide each affine piece into cells of mesh \(\eta\).
If a cell has left endpoint \(x_i\), length \(d_i\), left value \(y_i\), and drop
\(m_i=z_i d_i\le D\eta\), put
\[
 A_i=A(x_i),\qquad g_i=G(y_i),\qquad b_i=\min(z_i,2g_i).
\]
Insert at the left endpoint a potential drop \(m_i-b_id_i\), then
use an affine segment of slope \(-b_i\) reaching the prescribed
right value. The new potential stays in the cell's original range.
Writing it as \(\widetilde q\), the temperature \(u=Q(\widetilde q)\)
is piecewise \(H^1\), because \(Q'\) is bounded on \([\delta,Y]\).
There are finitely many jumps.

The frozen cell cost is exactly
\begin{equation}\label{eq:microjump}
 A_i\left\{\frac{d_ib_i^2}{4g_i}+m_i-b_id_i\right\}
                   =A_id_i\psi_{g_i}(z_i).
\end{equation}
On its affine part,
\(|u'|^2=b_i^2/(4G(\widetilde q))\),
\(|\widetilde q-y_i|\le m_i\), and \(b_i\le2g^*\).
The error between the actual and frozen Dirichlet terms is at most
\[
 \left(\frac{L_A\eta}{4g_*}
              +\frac{A_*L_Gm_i}{4g_*^2}\right)b_i^2d_i .
\]
The inserted jump has the exact weight \(A_i\).
The difference between the right side of \eqref{eq:microjump}
and the cost of the original affine cell of \(q\) is at most
\(L_A\eta m_i+A_*L_Gd_im_i\).
Since \(m_i\le D\eta\), \(\sum d_i\le R-1\), and
\(\sum m_i\le Y\), the sum of these errors is \(C_\delta\eta\)
for a finite constant depending on the fixed polygon, cutoff and
original parameters. Retained jumps have the same prescribed drop;
if jumps coincide, their potential differences add.
Thus
\[
 J(u)\le F(q)+C_\delta\eta .
\]
Choose first \(h\), then \(\delta\), and finally \(\eta\) to make
the total errors in this inequality, \eqref{eq:polygonalerror}, and
\eqref{eq:cutofferror} less than \(\epsilon\). No uniform bound on
\(1/G\) as \(\delta\downarrow0\) is used.
\end{proof}

\begin{proposition}\label{prop:elimination}
Every finite-jump temperature in \eqref{eq:J} can be replaced, without
increasing its cost, by an admissible temperature for
\((B_1,B_\ell)\) for some \(\ell\in[1,R]\). Its physical energy is
at most \(S J(u)\).
\end{proposition}

\begin{proof}
Take the first downward jump with positive right temperature. Denote
its left and right temperatures by \(b>c>0\), and its radius by \(s\).
On the whole tail after \(s\), replace the potential by
\(\lambda\Theta(u)\), where
\[
 0\le\lambda\le M:=\frac{\Theta(b)}{\Theta(c)}>1.
\]
Its temperature is \(g_\lambda(u)=Q(\lambda\Theta(u))\).
All tail temperatures are at most \(c\), so the inverse is evaluated
in \([0,\Theta(b)]\subset[0,Y]\).

For each fixed \(\lambda>0\), \(g_\lambda\) is Lipschitz on \([0,c]\).
Indeed, if \(\gamma>0\), the leading term of \(\Theta(t)\) at zero
is \(2\beta\gamma^3t^2\); if \(\gamma=0\), it is \(\beta t^5/5\).
Writing the respective order as \(m=2\) or \(m=5\), one obtains
\[
 g_\lambda(t)/t\longrightarrow\lambda^{1/m},\qquad
 g_\lambda'(t)=\frac{\lambda\Theta'(t)}
                         {\Theta'(g_\lambda(t))}
                         \longrightarrow\lambda^{1/m}.
\]
The derivative is bounded away from that endpoint as well.
For \(\lambda=0\) take the zero map. Thus every transformed piece
is \(H^1\), and both endpoint costs are finite.

For \(y=\Theta(u)>0\), the transformed Dirichlet density is
\[
 \frac{\lambda^2|y'|^2}{4G(\lambda y)}
                   =\frac{|y'|^2}{4y^2}H(\lambda y).
\]
At \(\lambda=0\) its value is zero. On \(y=0\), \(y'=0\) a.e.,
and the density is also set to zero there.
Lemma~\ref{lem:concavity} makes this expression concave in \(\lambda\).
Integration preserves that inequality, with finite nonnegative
terms. Every later jump cost, including the terminal one, is linear
in \(\lambda\). At the selected jump the cost is
\[
 A(s)\bigl[\Theta(b)-\lambda\Theta(c)\bigr],
\]
also affine. The total cost \(J_\lambda\) is concave on \([0,M]\);
therefore
\[
 J_1\ge(1-1/M)J_0+J_M/M.
\]
One of the two endpoint choices costs no more than the original profile.

At zero, the entire tail disappears and the selected jump becomes
the outer boundary. At \(M\), the selected jump disappears because
its right temperature becomes \(b\). The number of jumps with
positive right value strictly decreases; monotonicity prevents
new upward jumps. After finitely many repetitions only a terminal
jump to zero remains, at some \(\ell\in[1,R]\).
If the profile reaches zero continuously, this terminal jump has
size zero. Equal traces at the finitely many intervening pieces
cancel their distributional boundary terms. The resulting radial
temperature, extended by 1 in \(B_1\), is in \(H^1(B_\ell)\).
Its gradient and outer trace costs are
\[
 S\int_1^\ell A(r)|u'(r)|^2\dd r
              +S A(\ell)\Theta(u(\ell-))=S J(u).
\]
At \(\ell=1\) this is exactly the bare ball. In particular, the
comparison is now with the original physical energy.
\end{proof}

\begin{proof}[Proof of Theorem~\ref{thm:main}]
Fix any original admissible \(v\). Proposition~\ref{prop:lower}
gives \(w\) with \(S F(w)\le\mathcal E_\Omega(v)\).
For every \(\epsilon>0\), Proposition~\ref{prop:approximation}
gives \(J(u)\le F(w)+\epsilon\), and
Proposition~\ref{prop:elimination} gives an admissible concentric-ball competitor.
Lemma~\ref{lem:radial} ensures that the radial minimum exists.
Hence
\[
 \min_{1\le r\le R}E_\Theta(B_1,B_r)
        \le S J(u)\le S F(w)+S\epsilon
        \le\mathcal E_\Omega(v)+S\epsilon.
\]
Let \(\epsilon\downarrow0\), with the original parameters and
temperature fixed, then minimize over \(v\).
This proves the claimed inequality for the original measurable
cores and bounded Lipschitz outer sets, including the bare ball
and every smaller outer volume.
\end{proof}

\section*{AI-use disclosure}
OpenAI Codex was used to assist in developing and checking the argument,
checking algebra and references, and drafting this manuscript.
The manuscript is presented as an AI-assisted mathematical argument
for independent assessment.
\begin{thebibliography}{9}
\bibitem{AGNT}
V.~Amato, A.~Gentile, C.~Nitsch and C.~Trombetti,
\emph{On the gradient rearrangement of functions},
Mathematische Annalen (2024).
\href{https://doi.org/10.1007/s00208-024-02847-3}{doi:10.1007/s00208-024-02847-3}.

\bibitem{BNNT}
D.~Bucur, M.~Nahon, C.~Nitsch and C.~Trombetti,
\emph{Shape optimization of a thermal insulation problem},
Calculus of Variations and Partial Differential Equations
\textbf{61} (2022), article 186.
\href{https://doi.org/10.1007/s00526-022-02298-1}{doi:10.1007/s00526-022-02298-1}.

\bibitem{Gorny}
W.~G\'orny,
\emph{Functions of bounded variation and their applications},
author-hosted lecture notes, undated; version retrieved 15 October 2026.
\url{https://www.mat.univie.ac.at/~wgorny/Teaching/VSM_lecture_BV_functions.pdf}.

\bibitem{Hunter}
J.~K.~Hunter,
\emph{Notes on Partial Differential Equations}, Chapter~3:
Sobolev spaces, author-hosted lecture notes.
\url{https://www.math.ucdavis.edu/~hunter/pdes/ch3.pdf}.

\bibitem{LV}
G.~P.~Leonardi and G.~Vianello,
\emph{Free-boundary monotonicity for almost-minimizers of the relative perimeter},
Interfaces and Free Boundaries, online-first version (2025).
\href{https://doi.org/10.4171/IFB/544}{doi:10.4171/IFB/544}.
\end{thebibliography}

\end{document}
